\chapter{Conclusion}
\label{ch:conclusion}

\section{Summary of Contributions}
\label{sec:summary-contributions}

This project set out to answer a narrow but practical question: can a transformer-based model simplify public-facing health information without quietly losing the meaning that makes the information useful in the first place. The answer, on the evidence gathered here, is a qualified yes. A T5-small model fine-tuned on Cochrane-auto produces fluent, on-task simplifications where the same model without fine-tuning frequently does not, and it modestly outperforms a hand-written rule-based baseline on the metrics built specifically to measure simplification quality and meaning preservation. It does not yet outperform that same baseline on every measure, and a manual reading of individual outputs suggests why: at the training scale used for the confirmation run reported in \cref{ch:results-discussion}, the model has learned to make small, safe edits rather than the deeper restructuring that some sentences genuinely need.

Beyond the specific numbers, this project produced a working, documented, and reproducible pipeline covering dataset preparation, two baselines, a fine-tuning procedure, and a five-part evaluation combining readability formulas, overlap metrics, a simplification-specific metric, a semantic similarity measure, and a manual reading of outputs. Every design choice, including ones that reduced surface performance, such as declining to strip statistical detail in the heuristic baseline, was made deliberately and is documented in this thesis rather than left implicit.

\section{Limitations}
\label{sec:limitations}

The results in this thesis come from a confirmation run using 800 of the 7{,}075 available training pairs over three epochs, run on hardware without a dedicated GPU. This was a deliberate choice made to validate that the full pipeline behaves correctly end to end before committing to a longer run, and the thesis has been explicit throughout about which numbers come from this smaller run rather than presenting them as the final word on the model's ability. A full run over the complete training set for up to eight epochs is the natural next step, and \cref{sec:future-work} sets out what would change as a result.

The evaluation itself has narrower limitations too. BERTScore was computed using \texttt{distilbert-base-uncased} rather than the larger \texttt{roberta-large} model used in the original BERTScore paper, a choice made for practicality on the available hardware that may understate meaning-preservation scores slightly. The qualitative analysis in \cref{ch:results-discussion} examined three illustrative examples in detail rather than a large systematic sample, which is enough to identify a pattern but not enough to quantify how often that pattern occurs across the full test set. Finally, this project works entirely at sentence level, which sidesteps the document-level coherence questions raised in \cref{sec:variation-coherence} and leaves them open for future work.

\section{Ethical Considerations}
\label{sec:ethics}

This project uses only publicly available datasets, Cochrane-auto, ASSET, and WikiAuto, and involves no direct interaction with patients or clinical data. Even so, the application it points toward carries real ethical weight. A simplification model can produce text that reads as authoritative and clear while quietly dropping a limitation, a condition, or a note of uncertainty that a clinician would consider essential. It can also, in the opposite direction, make a claim sound more certain than the evidence actually supports, simply by replacing careful, hedged scientific language with a more direct sentence. Both risks were treated as central design concerns throughout this project rather than as an afternote, which is why meaning preservation was evaluated with the same seriousness as readability rather than as a secondary check.

This system is not positioned, and should not be read, as a medical decision-making tool. It is an assistive drafting aid at most, intended to sit alongside human review rather than replace it. The failure case examined in \cref{sec:qualitative-discussion}, where the model made no edit at all to a sentence that needed substantial rewriting, is arguably a safer failure mode than the alternative of confidently producing a fluent but inaccurate simplification, and that asymmetry is worth keeping in mind when judging what counts as an acceptable error rate for a system like this one.

\section{Future Work}
\label{sec:future-work}

Several extensions follow directly from the limitations above. The most immediate is completing a full training run across all 7{,}075 training pairs for up to eight epochs, which the pipeline already supports and which would give a fairer picture of what fine-tuning at this scale can achieve. Repeating the same pipeline with BART-base, noted throughout this thesis as an alternative architecture, would show whether the results reported here are specific to T5-small or hold more generally across compact transformer models. Re-running BERTScore with \texttt{roberta-large} would bring the meaning-preservation evaluation in line with the original metric's recommended configuration.

Beyond these direct extensions, two questions raised earlier in this thesis remain genuinely open. The dataset analysis in \cref{ch:dataset} suggested that a meaningful share of Cochrane lay summaries add explanatory content rather than simply compressing their source sentence, which points toward framing this task, at least in part, as constrained lay summarisation rather than sentence-level simplification in the narrow sense. A future iteration of this project could test that framing directly by training on paragraph-level or document-level pairs rather than isolated sentences. The coherence concerns raised in \cref{sec:variation-coherence} point in the same direction. A patient information document is read as a connected argument, not as a set of independently simplified sentences, and evaluating a health simplification system at document level would be a natural and worthwhile continuation of the work presented here.
